\documentclass[10pt,a4paper]{amsart}
\usepackage[margin=19mm]{geometry}
\usepackage{amsmath,amssymb,mathtools}
\usepackage[scr=rsfs,cal=euler]{mathalfa}
\usepackage{xfrac}
\usepackage[unicode,hidelinks,bookmarks=false]{hyperref}
\newcommand{\ie}{i.e.\ }
\newcommand{\cf}{cf.\ }
\newcommand{\resp}{respectively\ }
\newcommand{\Subsection}{\subsection}
\providecommand{\nobreakdash}{\nobreak\hbox{-}}
\setlength{\emergencystretch}{2em}
\multlinegap0pt
\usepackage{bm}
\usepackage{amssymb}
\usepackage{xcolor}
\usepackage{enumerate}
\newtheorem{lemma}{Lemma}
\renewcommand{\d}{\,{\rm d}}
\title{Mean Value Estimates for a Real-Exponent Analogue of Waring's Problem}
\author{Ataleshvara Bhargava}
\date{Source: arXiv:2608.18356v1 (CC BY 4.0)}
\begin{document}
\maketitle

\noindent\emph{Source and licence.} Mean Value Estimates for a Real-Exponent Analogue of Waring's Problem. Version arXiv:2608.18356v1 (CC BY 4.0), distributed by arXiv under the Creative Commons Attribution 4.0 International licence, http://creativecommons.org/licenses/by/4.0/, as recorded in the arXiv licence metadata for this exact version and shown on the abstract page https://arxiv.org/abs/2608.18356v1. Full licence notice: \texttt{licenses/arxiv-2608.18356-CC-BY-4.0.txt}.\par\smallskip

\noindent\textit{Editorial scope.} Counted unit: the article's lemma lem:meancasegen (source0.tex lines 508--511). The article's complete original proof of it is delivered with it (source0.tex lines 513--618, one \texttt{proof} environment of 106 lines). No mathematics is written, repaired or re-derived: the statement and the proof are the article's own lines, in the article's own notation, and no retained line is substituted or re-flowed.  Retained uncounted as the source-local context the proof needs: lines 232--240; lines 354--359.  Nothing was omitted from any retained span, so the package carries no omission record.  No retained line contains a citation, so the wrapper carries no bibliography and no bibliography entry.  No retained line contains a figure environment, an image-inclusion command or drawing code, so no image is delivered and the package declares no asset dependency.  Editorial formatting only, in addition: an AMS \texttt{amsart} preamble that restates the article's own declarations and macros used by the retained lines, namely the environments \texttt{lemma} on separate counters, together with the standard packages those lines need.  The retained environments print in this document as Lemma 1 in order of appearance, whereas the article prints the same items with the numbering of its own full document.  The complete native source of the article as distributed in the arXiv e-print for this exact version is bundled unchanged as \texttt{sources/207986/source0.tex}.
\begin{lemma} \label{lem:poulcount}
Define $\Delta = (2\delta)^{-1}$. Then for any $1\leq h \leq r$ and $n \in \mathbb{N}$ one has 
\[ V_{r,h}^n(\mathcal{S}, I_2,;\delta,m) \ll \delta \int_{-\Delta}^{\Delta} \int_0^1 |H_{\mathcal{S}}(\alpha,\boldsymbol{\beta}_n)H_2(\alpha,\boldsymbol{\beta}_n)^{2r-2h}| \d \boldsymbol{\beta}_n \d \alpha. \]
In particular, one has 
\[ V_{r,h}^n(I_1, I_2; \delta,m) \ll \delta \int_{-\Delta}^{\Delta} \int_0^1 |H_1(\alpha,\boldsymbol{\beta}_n)^{2h} H_2(\alpha,\boldsymbol{\beta}_n)^{2r-2h}| \d \boldsymbol{\beta}_n \d \alpha. \]
Moreover, it also holds that 
\[ V_{r,h}^n(I_1, I_2; \delta) \asymp \delta \int_{-\Delta}^{\Delta} \int_0^1 |H_1(\alpha,\boldsymbol{\beta}_n)^{2h} H_2(\alpha,\boldsymbol{\beta}_n)^{2r-2h}| \d \boldsymbol{\beta}_n \d \alpha. \]
All implicit constants are independent of $I_1$, $I_2$, $\mathcal{S}$, $\theta$ and $\delta$. 
\end{lemma}

\begin{lemma} \label{lem:meancase2}
Let $\theta> 3$ and $\kappa \geq 1$. Let $m_3 \in \mathbb{N}$ such that $(m_3,3)$ is a Diophantine pair. There exists an interval $I_0 \subseteq (X,2X]$ of the form $I_0 = I(\mathbf{q}^{(2)}) $ such that 
\begin{align*} \int_{-\kappa}^{\kappa} |g(\alpha)|^{2r} \d\alpha 
& \ll \kappa X^{\frac{4r}{3}+\frac{2m_3}{3}-1+\epsilon})\int_{-1}^{1} \int_0^1 |H(\alpha,\boldsymbol{\beta}_2;I_0)|^{2r-2m_3} \d \boldsymbol{\beta}_2 \d \alpha.  
\end{align*}
\end{lemma} 

\begin{lemma} \label{lem:meancasegen}
Let $\theta > 3$ be a non-integer and let $3 \leq n \leq \lfloor \theta \rfloor$ be an integer. Let $(m_k,k)$ be a sequence of Diophantine pairs for $3 \leq k \leq n$ and let $M_k = m_3+\cdots+m_k$. For $X, \kappa \geq 1$, there exists an interval $I_0 \subseteq (X,2X]$ of the form $I_0 = I(\mathbf{q}^{(n-1)})$ such that for every $\epsilon > 0$ one has 
\[ \int_{-\kappa}^{\kappa} |g(\alpha)|^{2r} \d\alpha \ll_{\epsilon} \kappa X^{2r(1-\frac{1}{n})+\frac{2M_n}{n}-\frac{n-1}{2}+\epsilon} \int_{-1}^{1} \int_{0}^{1} |H(\alpha,\boldsymbol{\beta}_{n-1};I_0 )|^{2r-2M_n} \d \boldsymbol{\beta}_{n-1} \d \alpha. \]
\end{lemma}

\begin{proof} We proceed by induction on $n$. By Lemma \ref{lem:meancase2} the case $n=3$ follows. Indeed, we have
\[ \int_{-\kappa}^{\kappa} |g(\alpha)|^{2r} \d\alpha \ll \kappa X^{2r(1-\frac{1}{3})+\frac{2m_3}{3}-\frac{3-1}{2}+\epsilon} \int_{-1}^1 \int_0^1 |H(\alpha,\boldsymbol{\beta}_{n-1};I_0 )|^{2r-2m_3} \d \boldsymbol{\beta}_{2} \d \alpha. \]
Therefore, we assume $\theta > 4$ and let $3 < n \leq \lfloor \theta \rfloor$ such that Lemma \ref{lem:meancasegen} holds for $n-1$. Then there exists an interval $I \subseteq (X,2X]$ such that $I = I(\mathbf{q}^{(n-2)})$ and 
\[ \int_{-\kappa}^{\kappa} |g(\alpha)|^{2r} \d\alpha \ll \kappa X^{2r(1-\frac{1}{n-1})+\frac{2M_{n-1}}{n-1}-\frac{n-2}{2}+\epsilon} \int_{-1}^{1} \int_{0}^{1} |H(\alpha,\boldsymbol{\beta}_{n-2};I )|^{2r-2M_{n-1}} \d \boldsymbol{\beta}_{n-2} \d \alpha. \]
Applying Lemma \ref{lem:poulcount} we have 
\begin{align*} 
\int_{-1}^{1} \int_{0}^{1} |H(\alpha,\boldsymbol{\beta}_{n-2};I )|^{2r-2M_{n-1}} \d \boldsymbol{\beta}_{n-2} \d \alpha \ll V_{r-M_{n-1}}^{n-2}(I;1/2).
\end{align*}
Let $r_0 = r-M_{n-1}$. The right hand side is the number of integer solutions $(x_1,\ldots, x_{2r_0}) \in I^{2r_0}$ to the system 
\begin{equation} \label{eq:degminustwo}
\begin{aligned}
|x_1^{\theta}+\cdots+x_{r_0}^{\theta}-x_{r_0+1}^{\theta}-\cdots-x_{r_0+1}^{\theta}| < 1/2, 
\\
\sum_{i=1}^{r_0} (x_i^j-x_{r_0+i}^j) = 0 \quad \text{ for } \quad 1 \leq j \leq n-2.
\end{aligned}
\end{equation}
As observed earlier, 
\[ I = I(\mathbf{q}^{(n-2)}) \subseteq \bigsqcup_{q=0}^{\lfloor X^{\frac{1}{n(n-1)}} \rfloor} I({\mathbf{q}^{(n-1)}|q}). \]
Define $I^*$ to be the interval that is the disjoint union on the right. Therefore, it follows that 
\[ V_{r_0}^{n-2}(I;1/2) \ll V_{r_0,m_n}^{n-2}(I, I^* ;1/2), \]
where we recall that the quantity on the right is the number of integer solutions $(x_1,\ldots, x_{2r_0}) $ to the system above with
\[ (x_1,\ldots, x_{m_n}, x_{r_0+1},\ldots,x_{r_0+m_n}) \in I^{2m_n} \] 
and $x_i \in I^*$ for all other $i$. 

Applying Lemma \ref{lem:poulcount}, we have 
\[ V_{r_0,m_n}^{n-2}(I, I^* ;1/2) \ll \int_{-1}^{1} \int_{0}^{1} |H(\alpha,\boldsymbol{\beta}_{n-2};I )|^{2m_n} |H(\alpha,\boldsymbol{\beta}_{n-2};I^* )|^{2r_0-2m_n} \d \boldsymbol{\beta}_{n-2} \d \alpha. \]
By definition of $I^*$, we use the Triangle Inequality to deduce that 
\begin{align*} |H(\alpha,\boldsymbol{\beta}_{n-2};I^* )|^{2r_0-2m_n} & \leq \Big( \sum_{q = 1}^{\lfloor X^{\frac{1}{n(n-1)}} \rfloor} |H(\alpha,\boldsymbol{\beta}_{n-2};I({\mathbf{q}^{(n-1)}|q}) )|  \Big)^{2r_0-2m_n}
\\ 
& \ll X^{\frac{2r_0}{n(n-1)}-\frac{2m_n}{n(n-1)}} |H(\alpha,\boldsymbol{\beta}_{n-2};I({\mathbf{q}^{(n-1)}|q_{n-1}}) )|^{2r_0-2m_n}
\end{align*}
for some $1 \leq q_{n-1} \leq \lfloor X^{\frac{1}{n(n-1)}} \rfloor$. Using this inequality and Lemma \ref{lem:poulcount} yet again, one has 
\begin{align*} 
\int_{-1}^{1} & \int_{0}^{1} |H(\alpha,\boldsymbol{\beta}_{n-2};I )|^{2m_n} |H(\alpha,\boldsymbol{\beta}_{n-2};I^* )|^{2r_0-2m_n} \d \boldsymbol{\beta}_{n-2} \d \alpha.
\\ 
& \ll X^{\frac{2r_0}{n(n-1)}-\frac{2m_n}{n(n-1)}} \int_{-1}^{1} \int_{0}^{1} |H(\alpha,\boldsymbol{\beta}_{n-2};I )|^{2m_n} |H(\alpha,\boldsymbol{\beta}_{n-2};I({\mathbf{q}^{(n-1)}|q_{n-1}}) )|^{2r_0-2m_n} \d \boldsymbol{\beta}_{n-2} \d \alpha 
\\ 
& \ll X^{\frac{2r-2M_{n-1}}{n(n-1)}-\frac{2m_n}{n(n-1)}} V_{r_0,m_n}^{n-2}(I,I({\mathbf{q}^{(n-1)}|q_{n-1}}),1/2).
\end{align*}
Combining this estimate with the inductive assumption we have 
\begin{align*} 
\int_{-\kappa}^{\kappa} |g(\alpha)|^{2r} \d\alpha \ll \kappa X^{2r(1-\frac{1}{n})+\frac{2M_{n-1}}{n}-\frac{n-2}{2}-\frac{2m_{n}}{n(n-1)}+\epsilon} V_{r_0,m_n}^{n-2}(I,I({\mathbf{q}^{(n-1)}|q_{n-1}}),1/2). 
\end{align*}
It therefore suffices to show for any $\epsilon > 0$ that 
\begin{align*} V_{r_0,m_n}^{n-2} & (I,I({\mathbf{q}^{(n-1)}|q_{n-1}}),1/2) 
\\ 
& \ll X^{\frac{2m_n}{n-1}-\frac{1}{2}+\epsilon} \int_{-1}^{1} \int_{0}^{1} |H(\alpha,\boldsymbol{\beta}_{n-1};I({\mathbf{q}^{(n-1)}|q_{n-1}}) )|^{2r-2m_n} \d \boldsymbol{\beta}_{n-1} \d \alpha. 
\end{align*}
By definition, $V_{r_0,m_n}^{n-2}(I,I({\mathbf{q}^{(n)}|q_{n-1}}),1/2)$ counts the integer solutions $(x_1,\ldots,x_{2r_0})$ to \eqref{eq:degminustwo} with 
\[ x_1,\ldots,x_{m_n},x_{r_0+1},\ldots,x_{r_0+m_n} \in I = (X_{\mathbf{q}^{(n-2)}}, X_{\mathbf{q}^{(n-2)}}+X^{1/(n-1)}]\] 
and 
\[x_i \in I({\mathbf{q}^{(n-1)}|q_{n-1}}) = (X_{\mathbf{q}^{(n-1)}|q_{n-1}}, X_{\mathbf{q}^{(n-1)}|q_{n-1}}+X^{1/n}] \] 
for all remaining $i$. Write $T_{n-1} = \lfloor X_{\mathbf{q}^{(n-1)}|q_{n-1}} \rfloor $ and $z_i = x_i-T_{n-1}$. Then 
\[ 1 \leq z_i \leq X^{1/n}+1 \quad \text{ for } \quad i = m_n+1,\ldots,r_0,r_0+m_n+1,\ldots,2r_0. \] Moreover, since 
\[ X_{\mathbf{q}^{(n-1)}|q_{n-1}} \in [X_{\mathbf{q}^{(n-2)}},X_{\mathbf{q}^{(n-2)}}+X^{1/(n-1)}],\] 
it follows that $|z_i| \leq X^{1/(n-1)}+1$ for $i = 1,\ldots,m_n,r_0+1,\ldots,r_0+m_n$. Taylor expansion of the top line in \eqref{eq:degminustwo} allows us to deduce then that 
\begin{align*} 
1/2 & > \Big|\sum_{j=1}^{n-2} \binom{\theta}{j} T_{n-1}^{\theta-j} \sum_{i=1}^{r_0} (z_i^j-z_{r_0+i}^j) + \binom{\theta}{n-1}T_{n-1}^{\theta-n+1} \sum_{i=1}^{r_0} (z_i^{n-1} -z_{r_0+i}^{n-1}) + O\Big(X^{\theta-n}\sum_{i=1}^{2r_0} |z_i|^n\Big) \Big| 
\\ 
& = \Big| \binom{\theta}{n-1}T_{n-1}^{\theta-n+1} \sum_{i=1}^{r_0} (z_i^{n-1} -z_{r_0+i}^{n-1}) + O(X^{\theta-n+\frac{n}{n-1}}) \Big|
\\ 
& = \Big| \binom{\theta}{n-1}T_{n-1}^{\theta-n+1} \sum_{i=1}^{m_n} (z_i^{n-1} -z_{r_0+i}^{n-1}) +O(X^{\theta-n+1+\frac{n-1}{n}}) + O(X^{\theta-n+\frac{n}{n-1}}) \Big|,
\end{align*}
where in the last equality we used the bound $|z_i|, |z_{r_0+i|} \ll X^{1/n}$ for $i \neq 1,\ldots,m_n$. Hence, 
\[ |z_1^{n-1}+\cdots+z_{m_n}^{n-1}-z_{r_0+1}^{n-1}-\cdots-z_{r_0+m_n}^{n-1}| \ll X^{\frac{n-1}{n}}. \] 
Since the system of equations described by the second line in \eqref{eq:degminustwo} is translation invariant in the sense that $(a_1,\ldots,a_{2r_0})$ is a solution if and only if $(a_1+b,\ldots, a_{2r_0}+b)$ is, we have that $(z_1,\ldots,z_{2r_0})$ is a solution to that system as well. Thus, 
\[ \sum_{i=1}^{r_0} (z_i^j-z_{r_0+i}^j) = 0 \quad \text{ for } \quad 1 \leq j \leq n-2, \]
and therefore 
\[ \sum_{i=1}^{m_n} (z_i^j-z_{r_0+i}^j) = -\sum_{i=m_n+1}^{r_0} (z_i^j-z_{r_0+i}^j) \ll X^{j/n} \quad \text{ for } \quad 1 \leq j \leq n-2. \]
Hence, there exist constants $C_1,\ldots, C_{n-1}$ such that 
\[ \Big| \sum_{i=1}^{m_n} (z_i^j - z_{r_0+i}^j ) \Big| \leq C_j X^{\frac{j}{n}} \quad \text{ for } \quad 1 \leq j \leq n-1. \]
Let $\mathcal{S}_{n}$ be the set of tuples $(x_1,\ldots,x_{r_0+m_n}) \in I^{2m_n}$ such that the variables $z_i=x_i-T_{n-1}$ are solutions to the above system. Since $(m_n,n)$ is a Diophantine pair, we have $|\mathcal{S}_{n}| \ll X^{\frac{2m_n}{n-1}-\frac{1}{2}+\epsilon}$. Here, note that we have used $Y = X^{\frac{1}{n-1}}$. Applying Lemma \ref{lem:poulcount}, it follows that 
\begin{align*} 
V_{r_0,m_n}^{n-2} & (I,I({\mathbf{q}^{(n-1)}|q_{n-1}}),1/2) 
\\ 
&\ll \int_{-1}^{1} \int_{0}^{1} |H_{\mathcal{S}_2}(\alpha, \boldsymbol{\beta}_{n-2})| |H(\alpha,\boldsymbol{\beta}_{n-2};I({\mathbf{q}^{(n-1)}|q_{n-1}}) )|^{2r_0-2m_n} \d \boldsymbol{\beta}_{n-2} \d \alpha 
\\ 
& \ll X^{\frac{2m_n}{n-1}-\frac{1}{2}+\epsilon} \int_{-1}^{1} \int_{0}^{1} |H(\alpha,\boldsymbol{\beta}_{n-2};I({\mathbf{q}^{(n-1)}|q_{n-1}}) )|^{2r_0-2m_n} \d \boldsymbol{\beta}_{n-2} \d \alpha. 
\\ 
& \ll X^{\frac{2m_n}{n-1}-\frac{1}{2}+\epsilon} V_{r_0-m_n}^{n-2}(I({\mathbf{q}^{(n-1)}|q_{n-1}});1/2). 
\end{align*}
Now we have $z_i = x_i-T_{n-1} \in [1,X^{1/n}+1]$ for each $i$, and again these variables satisfy \eqref{eq:degminustwo} with $r_0$ replaced with $r_1 = r_0-m_n=r-M_n$. Taylor expanding again, we therefore deduce that 
\begin{align*} 
1/2 & > \Big| \sum_{j=1}^{n-2} \binom{\theta}{j} T_{n-1}^{\theta-j} \sum_{i=1}^{r_1} (z_i^j-z_{r_1+i}^j) + \binom{\theta}{n-1}T_{n-1}^{\theta-n+1} \sum_{i=1}^{r_1} (z_i^{n-1} -z_{r_1+i}^{n-1}) + O\Big(X^{\theta-n}\sum_{i=1}^{2r_1} |z_i|^n\Big) \Big|
\\ 
& = \Big| \binom{\theta}{n-1}T_{n-1}^{\theta-n+1} \sum_{i=1}^{r_1} (z_i^{n-1} -z_{r_1+i}^{n-1}) + O(X^{\theta-n+1}) \Big|,
\end{align*}
and hence 
\[ \sum_{i=1}^{r_1} (z_i^{n-1} -z_{r_1+i}^{n-1}) \ll 1. \]
This shows, after applying Lemma \ref{lem:poulcount}, that 
\begin{align*}
V_{r_1}^{n-2}(I({\mathbf{q}^{(n-1)}|q_{n-1}});1/2) & = \sum_{m \ll 1} V_{r_1}^{n-1}(I({\mathbf{q}^{(n-1)}|q_{n-1}});1/2, m)
\\ 
& \ll \sum_{m \ll 1} \int_{-1}^{1} \int_{0}^{1} |H(\alpha,\boldsymbol{\beta}_{n-1};I({\mathbf{q}^{(n-1)}|q_{n-1}}) )|^{2r_1} \d \boldsymbol{\beta}_{n-1} \d \alpha
\\ 
& \ll \int_{-1}^{1} \int_{0}^{1} |H(\alpha,\boldsymbol{\beta}_{n-1};I({\mathbf{q}^{(n-1)}|q_{n-1}}) )|^{2r_1} \d \boldsymbol{\beta}_{n-1} \d \alpha. 
\end{align*}
Recalling that $r_1 = r_0-m_n = r-M_n$, it follows by combining the estimates that 
\begin{align*} 
\int_{-\kappa}^{\kappa} |g(\alpha)|^{2r} \d\alpha & \ll \kappa X^{2r(1-\frac{1}{n})+\frac{2M_{n-1}}{n}-\frac{n-2}{2}-\frac{2m_{n}}{n(n-1)}+\epsilon} V_{r_0,m_n}^{n-2}(I,I({\mathbf{q}^{(n-1)}|q_{n-1}}),1/2) \\ 
& \ll \kappa X^{2r(1-\frac{1}{n})+\frac{2M_{n-1}}{n}-\frac{n-1}{2}+\frac{2m_{n}}{n}+\epsilon} V_{r_0-m_n}^{n-2}(I({\mathbf{q}^{(n-1)}|q_{n-1}});1/2)
\\ 
& \ll \kappa X^{2r(1-\frac{1}{n})+\frac{2M_{n}}{n}-\frac{n-1}{2}+\epsilon} \int_{-1}^{1} \int_{0}^{1} |H(\alpha,\boldsymbol{\beta}_{n-1};I({\mathbf{q}^{(n-1)}|q_{n-1}}) )|^{2r-2M_n} \d \boldsymbol{\beta}_{n-1} \d \alpha, 
\end{align*}
which completes the proof. 
\end{proof}

\end{document}
